% ---------------------------------------------------------------------------
% Keel -- merged paper 1+2 (see docs/DECISIONS.md D-042 for the positioning).
%
% STATUS: DRAFT. Sections 4-7 report experiments that are implemented, tested
% and reproducible from this repository. Section 8 specifies the hierarchical
% estimator and its abstention rule but reports NO results, because Phase 4 is
% not built -- do not remove its \todo marker until it does. Every number below
% is regenerable; papers/paper1/README.md lists the command for each table.
%
% Not compiled in this environment (no TeX installation). Compile with:
%   latexmk -pdf main.tex
% ---------------------------------------------------------------------------
\documentclass[11pt,a4paper]{article}

\usepackage[utf8]{inputenc}
\usepackage[T1]{fontenc}
\usepackage{amsmath,amssymb}
\usepackage{booktabs}
\usepackage{graphicx}
\usepackage{natbib}
\usepackage{hyperref}
\usepackage[margin=1in]{geometry}
\usepackage{xcolor}

\newcommand{\todo}[1]{\textcolor{red}{\textbf{[TODO: #1]}}}
\newcommand{\corr}{\mathrm{corr}}
\newcommand{\CATE}{\tau}

\title{\textbf{When Does Targeting on Treatment Effects Pay?}\\
\large The correlation between effect and propensity governs it,\\
and small samples decide whether you should act at all}

\author{Prasham Jain\thanks{Correspondence: \texttt{prasham258@gmail.com}. Code,
data pipelines and every number in this paper:
\url{https://github.com/PrashamJ17/PBL-Proj}.}}

\date{Draft --- \today}

\begin{document}
\maketitle

\begin{abstract}
Uplift modelling --- targeting customers by their estimated treatment effect rather
than by their outcome propensity --- is widely recommended and inconsistently
useful. We show that whether it pays is governed by a single measurable quantity:
$\corr(\hat\CATE, \hat\pi)$, the correlation across customers between the estimated
treatment effect and the estimated outcome propensity. Across four real randomised
experiments and one simulator we find that when the two orderings coincide, the
outcome model wins, not because it estimates the right quantity but because it
solves an easier one; and that as the correlation goes negative, the advantage of
effect-based targeting grows by an order of magnitude. Subscription retention is the
adversarial case, with a negative correlation created by a mechanism absent from
advertising and promotion: the customers a churn model ranks highest include dormant
payers for whom being contacted is itself the reminder to cancel. We record a
prediction made before the data was obtained --- that retail promotion would fall
between advertising and retention --- and it landed at $+0.18$ as predicted, though
underpowered to test the consequence.

We then show that this diagnosis is of limited practical use to the businesses it
most concerns, because at their scale neither method is reliable. Holding the
evaluation set fixed and shrinking only the training set on a real 64{,}000-customer
experiment, the best effect-based method beats random targeting on only 75\% of
draws at $n=500$; one standard estimator manages 55\%. Mean performance at that size
looks respectable. A business gets one draw.

We argue the correct response is a decision rule that can \emph{abstain}, and that
abstention must be denominated in money rather than probability: we give a
discrete-time survival and lifetime-value layer whose absolute probabilities are
calibrated well enough to multiply by revenue, and show that ranking customers by
value at risk overlaps a churn-risk ranking by only 21\% at the top decile. The
simulator used throughout is released as an instrument, not a contribution: it is
the only setting in which individual-level ground-truth treatment effects exist, and
no public dataset contains the mechanism the negative-correlation regime requires.
\end{abstract}

\medskip
\noindent\textbf{Keywords:} uplift modelling; heterogeneous treatment effects;
customer retention; small-sample decision-making; survival analysis; abstention.

\section{Introduction}
\label{sec:intro}

A firm with a fixed retention budget must choose which customers to treat. The
default in practice is to rank customers by predicted churn probability and treat
the top slice. \citet{ascarza2018} established with two field experiments that this
is \emph{ineffective} --- the highest-risk customers are not the most responsive ---
and that targeting on sensitivity to the intervention outperforms targeting on risk
by up to 6.8 percentage points. That finding is not in dispute here and we do not
restate it as our own.

Three questions remain open, and they are the ones a practitioner actually faces.

\paragraph{First: when is effect-based targeting worth its cost?} The uplift
literature reports wins on some datasets and not others
\citep{gutierrez2017,devriendt2018}, usually without a diagnosis. Estimating
$\CATE(x) = \mathbb{E}[Y(1)-Y(0) \mid X=x]$ is a harder statistical problem than
estimating $\pi(x) = \mathbb{E}[Y \mid X=x]$: it is a difference of two conditional
means, and its variance is the sum of theirs. If the two induce nearly the same
ranking, paying that variance premium buys nothing. Section~\ref{sec:corr} makes
this precise and measurable.

\paragraph{Second: can risk-based targeting be actively harmful, not merely
ineffective?} It can, but only under a condition that does not hold in most settings
where uplift modelling is studied. Section~\ref{sec:corr} states the condition and
Section~\ref{sec:instrument} describes the mechanism that produces it in
subscription retention.

\paragraph{Third: at the sample sizes of the firms that need this most, is any of it
reliable?} Published uplift results come from datasets of $10^4$ to $10^7$ customers.
The subscription businesses where retention economics are most acute have hundreds.
Section~\ref{sec:smalln} measures what happens in that regime and finds that the
honest answer is often \emph{no}, which motivates Section~\ref{sec:abstain}.

\paragraph{Contributions.}
\begin{enumerate}
  \item \textbf{A governing quantity.} $\corr(\hat\CATE,\hat\pi)$, measured on the
        $T$-learner so that both quantities live on comparable scales, orders five
        settings by how much effect-based targeting is worth
        (Table~\ref{tab:corr}). We give the mechanism, not only the correlation.
  \item \textbf{An out-of-sample confirmation.} The principle was used to predict
        where a fourth dataset would fall \emph{before it was obtained}. It landed
        as predicted. We also state plainly what that experiment failed to confirm.
  \item \textbf{Reliability, not expectation, at small $n$.} We report the
        proportion of draws on which a method beats random, paired on the same
        split, rather than a mean across draws.
  \item \textbf{A money-denominated decision layer.} A discrete-time competing-risks
        hazard model whose absolute survival probabilities are calibrated well
        enough to multiply by revenue, with an exact decomposition of lost value by
        cause of loss.
  \item \textbf{An instrument.} A calibrated subscription simulator with exact
        individual-level potential outcomes, released so that CATE estimators can be
        scored against per-customer truth rather than against aggregate proxies.
\end{enumerate}

\section{Related work}
\label{sec:related}

\paragraph{Uplift modelling.} The estimation problem has a long applied literature
\citep{radcliffe2011,rzepakowski2012} and a more recent causal-machine-learning one
\citep{athey2016,kunzel2019,chernozhukov2018}. Evaluation is usually by the Qini or
uplift curve. We report Qini for comparability but argue in
Section~\ref{sec:smalln} that it answers the wrong question for a firm with one
sample and a budget: it presumes you rank and treat a top slice, with no option to
decline.

\paragraph{Targeting for retention.} \citet{ascarza2018} is the closest prior work
and the reason our contribution is framed as \emph{when} risk-based targeting
becomes harmful rather than \emph{that} it is ineffective. \citet{hitsch2018} study
targeting under heterogeneous effects in marketing more broadly.

\paragraph{Cost-sensitive causal decisions.} \citet{verbeke2023} and
\citet{devriendt2021} formalise treatment decisions under explicit costs and
benefits, which is the same object as our Section~\ref{sec:value}. Our addition is
the survival and lifetime-value layer that supplies the benefit term, and the
requirement that its probabilities be calibrated rather than merely ranked.

\paragraph{Survival modelling of churn.} Cox regression \citep{cox1972} dominates,
with random survival forests \citep{ishwaran2008} and neural extensions
\citep{katzman2018} as the standard comparators. We use the person-period
discrete-time formulation \citep{singer1993}, and evaluate with proper scoring
rules \citep{graf1999} and distribution calibration \citep{haider2020} rather than
concordance alone.

\paragraph{Lifetime value.} We treat only the contractual case, where churn is
observed. The non-contractual (BTYD) family \citep{fader2005} is a different
estimand and is out of scope.

\section{The decision problem}
\label{sec:problem}

Let each customer $i$ have covariates $X_i$, a binary treatment $W_i$, and potential
outcomes $Y_i(0), Y_i(1)$ where $Y=1$ denotes churn within a horizon $H$. Write
\[
  \CATE(x) = \mathbb{E}[Y(1) - Y(0) \mid X = x], \qquad
  \pi(x)   = \mathbb{E}[Y \mid X = x].
\]
Under our sign convention $\CATE < 0$ is a treatment that \emph{helps}. Customers
with $\CATE > 0$ are \emph{sleeping dogs}: contacting them causes churn.

A policy $\rho: \mathcal{X} \to \{0,1\}$ under unit cost $c$ and customer value $V_i$
earns
\begin{equation}
  \label{eq:value}
  \mathcal{V}(\rho) = \sum_i \rho(X_i)\bigl(-\CATE(X_i) V_i - c\bigr),
\end{equation}
measured against \emph{doing nothing}. This baseline is not incidental. The metric
reported in most vendor case studies is the save rate among the treated, which is
contaminated by selection: target customers who were going to stay and the save rate
is superb while Equation~\eqref{eq:value} does not move.

Three properties of Equation~\eqref{eq:value} drive everything that follows. It
depends on $\CATE$, not $\pi$. It is denominated in money, so $V_i$ matters as much
as $\CATE_i$ --- the subject of Section~\ref{sec:value}. And it admits
$\rho \equiv 0$, so a policy that declines to act is a legitimate competitor ---
the subject of Section~\ref{sec:abstain}.

\section{The instrument}
\label{sec:instrument}

Individual treatment effects are never observed. Field experiments identify average
effects; public uplift datasets identify subgroup effects by randomisation. Neither
can score a per-customer CATE estimate against truth. We therefore build a simulator
in which both potential outcomes are computed exactly, and we are explicit that this
is an \emph{instrument} rather than a finding.

\paragraph{Structure.} A monthly discrete-time hazard drives voluntary churn;
involuntary churn (payment failure) is a separate process and the two are never
summed. Latent customer attributes --- engagement, habit strength, price
sensitivity, attention --- are drawn from a Gaussian copula. Observable features
mimic what a billing plus product-analytics integration actually provides. Latents
are returned separately from the observable panel and a test enforces that they
never enter it.

\paragraph{The mechanism that matters.} Latent \emph{attention} is correlated with
latent \emph{engagement}. Low engagement is the strongest observable predictor of
churn, so a churn model ranks low-engagement customers highest. Low attention marks
a dormant payer who has half-forgotten the charge, and whom contacting \emph{causes}
to cancel. Because the two latents are correlated, the top of a churn ranking is
dense in sleeping dogs. This is the entire negative-correlation regime, expressed as
a generative mechanism rather than as a fitted relationship.

\paragraph{Calibration, and one deliberate handicap.} Generative parameters are
solved for, not tuned by eye: the hazard intercept is obtained by bisection against
a target monthly churn rate. The simulator reproduces published SMB-SaaS aggregates
--- 4.5\% monthly voluntary churn, a 30\% involuntary share, 23\% two-year
retention, and an early-to-late hazard ratio of 2.0 --- and is stable across seeds.

One configuration choice runs \emph{against} our own thesis and is worth stating.
An earlier setting produced 30\% sleeping dogs and a mean treatment effect that was
\emph{harmful} ($\bar\CATE = +0.0044$). Under it the headline result is trivially
true: a blanket campaign raises churn, so of course targeting it badly loses money.
We rejected that configuration in favour of one with $\approx17\%$ sleeping dogs and
$\bar\CATE = -0.010$, so that a blanket campaign \emph{helps} on average and any
money lost is attributable to \emph{targeting} alone. This is the harder claim. A
build-time gate fails if $\bar\CATE$ ever becomes positive.

\paragraph{The kill test.} We train a real churn model --- gradient-boosted trees on
observable features only, with a strictly temporal split, reaching holdout
AUC $\approx 0.70$ --- and compare targeting policies at a 20\% budget over six
seeds (Table~\ref{tab:killtest}). Ranking by churn score loses money and loses
\emph{more than random selection}, unanimously across seeds. The mechanism is
visible directly: sleeping dogs are 48\% of the top predicted-risk decile and 2\% of
the bottom.

\begin{table}[t]
\centering
\begin{tabular}{lrrr}
\toprule
Policy & Treated & Value vs.\ do-nothing & Sleeping dogs treated \\
\midrule
Random (20\%)                & 718 & $-17{,}035$ & 9 \\
Churn-score top 20\%         & 718 & $-22{,}123$ & 19 \\
Effect-based top 20\%        & 718 & $+5{,}877$  & 0 \\
\textbf{Effect-based, with abstention} & \textbf{209} & $\mathbf{+8{,}610}$ & \textbf{0} \\
\bottomrule
\end{tabular}
\caption{Simulator, one representative seed at a 20\% budget; the ordering is
unanimous over six seeds (mean: churn-score $-26{,}903$ vs.\ random $-16{,}949$).
The last row treats 29\% as many customers as the third and earns 46\% more. Note
that a \emph{more accurate} churn model makes the second row worse, because accuracy
at detecting disengagement is precision at finding sleeping dogs.}
\label{tab:killtest}
\end{table}

\paragraph{A methodological note on leakage.} Every feature in this project is
computed point-in-time: each fact carries both when it occurred and when it became
\emph{knowable}, and features filter on the latter. The cost of not doing so is
measured rather than asserted. The same features, computed three ways, give holdout
AUC $0.603$ (correct), $0.613$ (ignoring reporting lag), and $\mathbf{0.954}$ (no
temporal filter). The 0.35 gap is not performance; it is the amount by which a
backtest would overstate a model before it reached production. We report it because
several published churn results are, on inspection, in the third column.

\section{What governs whether effect-based targeting pays}
\label{sec:corr}

\subsection{An apparent contradiction}

On the Hillstrom email experiment \citep{hillstrom2008}, effect-based targeting
clearly beats outcome-based targeting (Qini 168 vs.\ 36). On the Criteo advertising
experiment \citep{diemert2018} it clearly loses, at \emph{every} training size
(3{,}617 vs.\ 2{,}811 incremental visits at $n=500$). Neither result is a fluke and
they cannot both be evidence for a general recommendation.

\subsection{The reconciliation}

Estimate both quantities on held-out data and measure their correlation across
customers. To keep the comparison meaningful we measure it on the $T$-learner
specifically, whose output is a difference of two probabilities and therefore lives
on the treatment-effect scale; class-transform scores are scale-dependent under
unbalanced assignment and produce spurious correlations near unity.

\begin{table}[t]
\centering
\begin{tabular}{lrrr}
\toprule
Setting & $\corr(\hat\CATE, \hat\pi)$ & Advantage of effect-based & Predicted harmed \\
\midrule
Hillstrom (mens)     & $+0.63$ & $+4.7\%$   & 0.2\% \\
Criteo (advertising) & $+0.61$ & $+0.6\%$   & 19.2\% \\
Lenta (retail SMS)   & $+0.18$ & $+10.4\%$  & 24.9\% \\
Hillstrom (womens)   & $+0.07$ & $+3.9\%$   & 10.2\% \\
\textbf{Subscription retention (simulator)} & $\mathbf{-0.19}$ & $\mathbf{+106.8\%}$ & 25.7\% \\
\bottomrule
\end{tabular}
\caption{The governing quantity. Four settings is a contrast, not a fitted
relationship, and the ordering among the three positive-correlation rows is within
noise. The signal is the order-of-magnitude gap at negative correlation.}
\label{tab:corr}
\end{table}

An outcome model ranks by \emph{likelihood of responding}; an effect model ranks by
\emph{how much treatment changes the response}. When those orderings coincide, the
outcome model wins --- not because it measures the right thing, but because it
solves an easier estimation problem, and at small $n$ that variance advantage
dominates. This reframes the recommendation. It is not that uplift modelling is
better; that is false in advertising and Table~\ref{tab:corr} shows it. It is that
\textbf{retention has an adversarial structure that advertising does not}.

\subsection{The share of harmed customers is not the right diagnostic}

A natural hypothesis is that effect-based targeting pays when sleeping dogs exist.
Table~\ref{tab:corr} refutes it: Criteo has 19.2\% predicted-harmed customers, more
than Hillstrom-womens' 10.2\%, and effect-based targeting adds essentially nothing
there. Lenta has 24.9\%, comparable to the retention setting's 25.7\%, and also buys
little. What matters is not how many customers are harmed but \emph{whether the
harmed ones sit where the outcome model ranks highest} --- which is what a negative
correlation measures.

\subsection{A prediction, recorded before the data was obtained}

A relationship inferred from three points is a description. To make it a claim we
recorded a prediction in the repository before downloading a fourth dataset: retail
SMS promotion (Lenta, 687{,}029 customers; \citealp{sklift}) should fall
\emph{between} advertising
($+0.61$) and subscription retention ($-0.19$).

\textbf{Result: $+0.177$.} As predicted.

We separate what this confirms from what it does not. \emph{Confirmed:} the
correlation lands where predicted. \emph{Not confirmed:} whether the consequence
follows. Lenta's nominal $+10.4\%$ advantage has a bootstrap interval of
$[359, 1601]$ against $[309, 1477]$ for the outcome model and $[343, 1007]$ for
random --- nothing is distinguishable from anything. The dataset is underpowered
rather than noisy: its average treatment effect is $+0.0075$ on a 10.3\% base, a
7.4\% lift against Hillstrom's 42.6\%, and its small-$n$ sweep is flat from $n=500$
to $n=20{,}000$ because there is barely a signal to learn. This is one confirming
observation, not two.

\section{Reliability at the sample sizes that matter}
\label{sec:smalln}

Every public uplift dataset is large by small-business standards. The firms for whom
retention economics are most acute have hundreds of customers. We therefore hold the
\emph{evaluation} set fixed on Hillstrom and shrink only the training set, so any
change is attributable to estimation difficulty rather than evaluation noise.

The headline metric is deliberately not a mean. A mean across hypothetical draws is
not what a business experiences; it gets one draw, and a method with a good average
and a wide spread is a gamble. We therefore report the \textbf{proportion of seeds on
which a method beats random targeting, paired on the same split}.

\begin{table}[t]
\centering
\begin{tabular}{lrrrr}
\toprule
Training $n$ & 500 & 1{,}000 & 2{,}000 & 5{,}000+ \\
\midrule
Best effect-based method beats random & \textbf{75\%} & 90\% & 100\% & 100\% \\
Class-transform estimator beats random & \textbf{55\%} & --- & --- & --- \\
\bottomrule
\end{tabular}
\caption{Hillstrom, evaluation set held fixed. At $n=500$ one standard estimator is
a coin flip while its mean performance looks respectable.}
\label{tab:smalln}
\end{table}

Two consequences. First, published uplift results are reported in a regime almost no
prospective user occupies, and the gap is not a matter of degree: at $n=500$ the
question is not which method is better but whether \emph{any} of them is doing
anything. Second, the failure is invisible under the reporting conventions of the
field. Reporting the mean at $n=500$ would have shown a respectable number.

We find the same wall in the survival layer of Section~\ref{sec:value}: below
roughly 250 customers, no covariate model reliably beats the marginal survival curve
on a proper scoring rule. Different task, different metric, same conclusion.

\section{From probability to money}
\label{sec:value}

Equation~\eqref{eq:value} multiplies $\CATE_i$ by $V_i$. An identical treatment
effect on a \$12 customer and a \$1{,}200 customer are not comparable acts, and no
amount of churn-model accuracy distinguishes them. Supplying $V_i$ requires absolute
survival probabilities, not a ranking.

\paragraph{Discrete-time competing risks.} We expand each customer into one row per
period at risk and fit the hazard $h_k(t \mid x)$ for each cause $k$ separately;
survival is $\prod_{s \le t}\bigl(1 - \sum_k h_k(s)\bigr)$ and the cause-specific
cumulative incidence is $\sum_{s\le t} h_k(s) S(s-1)$. Causes are combined only at
the survival level and never at the label level, so voluntary churn and payment
failure keep separate models, separate coefficients and separate interpretations.
Three properties motivate this over Cox: time-varying covariates are native rather
than requiring a start-stop construction (which is where point-in-time discipline
usually breaks); the baseline hazard is estimated rather than profiled out as a
nuisance, which is what absolute calibration requires; and subscription billing is
genuinely discrete, so tied event times are simultaneity rather than a tie problem.

\paragraph{Calibration is the point, not concordance.} We evaluate on the Telco
customer-churn dataset \citep{telco} (7{,}032 real subscription customers), excluding
\texttt{TotalCharges}, which correlates with tenure at $r = 0.83$ and is therefore a
direct encoding of the duration being predicted --- an exclusion that published
Telco survival analyses routinely omit.

\begin{table}[t]
\centering
\begin{tabular}{lrrrr}
\toprule
Model & C-index & IBS $\downarrow$ & Calib.\ slope & Wins vs.\ ours \\
\midrule
\textbf{Discrete-time hazard (ours)} & 0.8650 & \textbf{0.0824} & 1.02 & --- \\
DeepSurv                             & 0.8661 & 0.0825 & 0.98 & 6/10 \\
Cox proportional hazards             & 0.8565 & 0.0914 & 1.18 & 0/10 \\
Random survival forest               & 0.8461 & 0.0964 & 1.18 & 0/10 \\
Kaplan--Meier (no covariates)        & ---    & 0.1823 & ---  & 0/10 \\
\bottomrule
\end{tabular}
\caption{Telco, mean over ten resplits. IBS is the integrated Brier score, a proper
scoring rule. We beat Cox and the random survival forest on 10/10 resplits and
\textbf{tie} DeepSurv --- $0.0824$ against $0.0825$ is noise and we report it as a
tie. On GBSG2, a continuous-time medical benchmark those methods were developed on,
we lose; we recorded that expectation before running it.}
\label{tab:survival}
\end{table}

Kaplan--Meier is in Table~\ref{tab:survival} for a reason. On calibration error
alone it \emph{beats} Cox on this dataset, because a covariate-free prediction is
trivially calibrated. Reporting calibration error as a headline would make
``predict the average for everyone'' the winning strategy. The integrated Brier
score penalises both miscalibration and lack of discrimination, and under it
Kaplan--Meier is not competitive at all.

\paragraph{Lifetime value, and a refusal.} With $S(t)$ in hand,
$\mathrm{CLV} = \sum_{t=1}^{H} m\, S(t)\,(1+d)^{-t}$ using contribution margin $m$
and a monthly discount $d$. We cap $H$ at the observed support and raise an error
otherwise. The standard $\mathrm{ARPU}/\text{churn rate}$ formula is an infinite
geometric sum that assumes a constant hazard forever; retention hazards decline with
tenure, so that extrapolation is biased rather than merely uncertain, and the
survival model has no evidence at all beyond its training window.

\paragraph{The result worth stating.} Ranking simulated customers by churn
probability and by value at risk ($\mathrm{CLV} \times$ hazard), \textbf{the two top
deciles overlap by 21\%}. A churn score points at the wrong four-fifths of the
money, before any causal argument is made. Separately, the value lost relative to
perfect retention decomposes \emph{exactly} into 72.5\% voluntary and 27.5\%
involuntary, with no residual --- a split a blended churn figure cannot express, and
one with entirely different remedies.

\section{Abstention under posterior uncertainty}
\label{sec:abstain}

\todo{Phase 4. This section specifies the estimator and the decision rule; it
reports no results, because the estimator is not yet built. Nothing in this section
should be read as demonstrated.}

Sections~\ref{sec:corr} and \ref{sec:smalln} together pose a problem a practitioner
cannot solve by choosing a better estimator. They cannot know in advance which
correlation regime they occupy, choosing wrongly is expensive in both directions,
and at their sample size the estimate that would tell them is itself unreliable.

The response we propose is a policy that quantifies its own uncertainty and declines
to act when that uncertainty is too wide relative to the stakes. Concretely, with a
posterior $p(\CATE_i \mid \mathcal{D})$ from a hierarchical model that partially
pools across customers (and, in a multi-tenant setting, across firms), treat $i$
only when
\begin{equation}
  \label{eq:abstain}
  \Pr\bigl(-\CATE_i V_i > c \;\big|\; \mathcal{D}\bigr) > 1 - \alpha ,
\end{equation}
which reduces to a point-estimate threshold as the posterior concentrates and to
$\rho \equiv 0$ when it does not. Partial pooling is the mechanism by which the
small-$n$ regime is survivable at all: it shrinks individual estimates toward a
population effect whose sign is far better identified than any individual's.

Table~\ref{tab:killtest} contains a crude precursor --- a threshold rule on an
oracle effect estimate treating 209 customers and earning more than a ranking rule
treating 718 --- but that uses ground-truth effects and is therefore an upper bound
on what an estimated rule can achieve, not evidence for one. The evaluation this
section requires is: does Equation~\eqref{eq:abstain}, with estimated posteriors,
beat both ranking and doing nothing on realised money at $n \in [250, 2000]$, on
what proportion of draws?

\section{Limitations}
\label{sec:limitations}

\paragraph{One simulator carries the negative-correlation regime.} Every setting in
Table~\ref{tab:corr} with $\corr < 0$ is ours. We argue the mechanism is real and
the correlation is a consequence of it rather than an input --- the attention--%
engagement correlation is a generative parameter and the correlation of
\emph{estimated} quantities is measured downstream --- but a public retention
experiment with the mechanism present would be far stronger evidence, and none
exists that we are aware of.

\paragraph{Four points are a contrast, not a relationship.} We do not fit a curve
through Table~\ref{tab:corr} and do not claim one. The ordering among the three
positive rows is within noise.

\paragraph{The Lenta confirmation is partial.} It confirms the axis and is
underpowered to test the consequence, as stated in Section~\ref{sec:corr}.

\paragraph{The forward window is less stochastic than the historical period.} The
simulator's exact-$\CATE$ regime evolves state along its expected trajectory, which
slightly understates outcome variance. Acceptable because the estimand is the
treatment effect rather than the outcome distribution, but it is a real
approximation.

\paragraph{The value layer is honest about probabilities, not about margins.}
Contribution margin and discount rate are inputs a firm supplies. Every monetary
figure in Section~\ref{sec:value} is conditional on them and they are stated
wherever a number appears.

\paragraph{The central proposal is unevaluated.} Section~\ref{sec:abstain} is a
specification. Until it is built and tested at small $n$, this paper is a diagnosis
with a proposed treatment, and should be read as one.

\section{Conclusion}

Whether to target on treatment effects is not a methodological preference but an
empirical question about a specific dataset, and $\corr(\hat\CATE, \hat\pi)$ answers
it. Where the effect ordering and the propensity ordering coincide --- advertising,
promotion --- the simpler estimator wins on variance grounds and the sophistication
is wasted. Where they diverge, as they do in subscription retention because the
customers a churn model ranks highest include those an intervention harms, the gap
is an order of magnitude and risk-based targeting destroys value rather than merely
failing to add it.

The practical difficulty is that the firms in the adversarial regime are also the
smallest, and at their scale neither estimator is reliable on a single draw. We take
that to mean the right object is not a better ranker but a decision rule that can
decline, denominated in money rather than probability. We have built the money and
shown that it disagrees sharply with risk. The decision rule itself remains to be
demonstrated.

\subsection*{Reproducibility}

Every number in this paper is regenerated by a command in the repository, listed in
\texttt{papers/paper1/README.md}. The simulator, the point-in-time feature store, the
benchmark harness, the survival and value layers, and 277 tests are released under
the repository's licence. Predictions made in advance of experiments are recorded in
the source files that run them, in commits that precede the runs.

\bibliographystyle{plainnat}
\bibliography{refs}

\end{document}
